\documentclass[11pt]{article}

\usepackage[a4paper,top=0.3in,bottom=0.3in,left=0.3in,right=0.3in]{geometry}
\usepackage{amsmath,amssymb,amsfonts,amsthm,mathtools}
\usepackage{bm}
\usepackage{enumitem}
\usepackage{booktabs}
\usepackage{algorithm}
\usepackage{algorithmic}
\usepackage{hyperref}
\usepackage{microtype}
\usepackage{cleveref}
\usepackage[normalem]{ulem}

% ── Custom commands ──
\newcommand{\R}{\mathbb{R}}
\newcommand{\N}{\mathbb{N}}
\newcommand{\Th}{\mathcal{T}_h}
\newcommand{\V}{V}
\newcommand{\Vh}{V_h}
\newcommand{\Vp}{V_h^P}
\newcommand{\Vhb}{V_h^b}
\newcommand{\dom}{\Omega}
\newcommand{\hh}{\mathbf{h}}
\newcommand{\Pe}{\operatorname{Pe}}
\newcommand{\calL}{\mathcal{L}}
\newcommand{\calT}{\mathcal{T}}
\newcommand{\calB}{\mathcal{B}}
\newcommand{\dd}{\,\mathrm{d}}
\newcommand{\supp}{\operatorname{supp}}
\newcommand{\tr}{\operatorname{tr}}

% ── Theorem environments ──
\newtheorem{definition}{Definition}[section]
\newtheorem{theorem}{Theorem}[section]
\newtheorem{lemma}[theorem]{Lemma}
\newtheorem{proposition}[theorem]{Proposition}
\newtheorem{corollary}[theorem]{Corollary}
\theoremstyle{remark}
\newtheorem{remark}{Remark}[section]

% ── Title ──
\title{Learning Residual-Free Bubbles with Kolmogorov--Arnold Networks for Singularly Perturbed Elliptic Problems}
\author{Juan Felipe Pacazuca Santiago}
\date{\today}

\begin{document}
\maketitle

% ════════════════════════════════════════════════════════════════════════
\begin{abstract}
We present a method for learning residual-free bubble (RFB) functions in one-dimensional advection--diffusion--reaction problems using Kolmogorov--Arnold Networks (KANs). The standard Galerkin finite element approximation with piecewise linear elements suffers from severe accuracy degradation when the reaction or advection coefficients dominate the diffusion---the so-called singularly perturbed regime. Residual-free bubbles cure this deficiency by enriching the coarse finite element space with local functions that exactly cancel the element residual. However, computing these bubbles requires solving a local boundary-value problem on every element for every new physical configuration.

We propose to replace the exact local solves with a KAN-based surrogate that maps the two local dimensionless parameters---the element P\'{e}clet number $\Pe_K$ and the reaction number $\rho_K$---to the bubble shape on the reference element. The network architecture is a true KAN: each edge carries a learnable univariate function composed of a SiLU base and a B-spline correction, rather than a fixed activation applied to a linear combination. The trained model is mesh-independent by construction, since the input parameters are non-dimensional. Bubble degrees of freedom are eliminated statically, so the final global system has the same dimension as the classical Galerkin method.

We provide a complete mathematical framework: the weak formulation in $H^1_0(\Omega)$, the Galerkin approximation, the abstract RFB theory, static condensation, and a rigorous error estimate that decomposes the total error into a finite element approximation error and a bubble approximation error. Numerical experiments confirm that the KAN-RFB method achieves accuracy comparable to the exact RFB across multiple mesh refinements, demonstrating true mesh independence.
\end{abstract}

% ════════════════════════════════════════════════════════════════════════
\section{Introduction}
\label{sec:introduction}

\subsection{The model problem}

We consider the one-dimensional reaction--advection--diffusion equation with possibly non-homogeneous Dirichlet boundary conditions:
\begin{equation}\label{eq:PDE}
\begin{cases}
\displaystyle -\frac{d}{dx}\!\left(\kappa\frac{du}{dx}\right) + \beta\frac{du}{dx} + \sigma\,u = f & \text{in } \dom = (0,\ell),\\[4pt]
u(0) = g_0,\quad u(\ell) = g_\ell,
\end{cases}
\end{equation}
where $\kappa > 0$ is the diffusivity (or thermal conductivity, permeability, etc.), $\beta \in L^\infty(\dom)$ is the advective velocity, $\sigma \in L^\infty(\dom)$ is the reaction (or absorption) coefficient, $f \in L^2(\dom)$ is the source term, and $g_0, g_\ell \in \R$ are prescribed boundary data.

This equation arises in a wide range of applications:
\begin{itemize}[nosep]
    \item \textbf{Heat transfer:} $\kappa$ is thermal conductivity, $\beta$ represents convective transport, $\sigma$ models radiative cooling.
    \item \textbf{Groundwater flow:} $\kappa$ is hydraulic conductivity, $\beta$ is the Darcy velocity, $\sigma$ accounts for chemical decay.
    \item \textbf{Pollutant transport:} the equation models advection--dispersion with first-order degradation in porous media.
    \item \textbf{Semiconductor device simulation:} $\kappa$ is the diffusion coefficient of carriers, $\beta$ is the drift velocity under an electric field.
\end{itemize}

The essential feature that makes \eqref{eq:PDE} mathematically and numerically challenging is the interplay of three competing mechanisms---diffusion, advection, and reaction---characterized by two dimensionless ratios. On an element of size $h$, these are the local P\'{e}clet number and the reaction number:
\begin{equation}\label{eq:dimensionless}
\Pe_K := \frac{|\beta|\,h_K}{2\kappa}, \qquad \rho_K := \frac{\sigma\,h_K^2}{\kappa}.
\end{equation}
When either $\Pe_K \gg 1$ (advection-dominated) or $\rho_K \gg 1$ (reaction-dominated), the problem is called \emph{singularly perturbed} and the solution develops thin boundary or interior layers that the coarse mesh cannot resolve.

\subsection{Sobolev spaces and weak formulation}

Let $\dom = (0,\ell)$ and define the Sobolev space $H^1(\dom)$ as the space of functions $v \in L^2(\dom)$ with weak derivative $v' \in L^2(\dom)$, equipped with the norm
\begin{equation}
\|v\|_{H^1(\dom)} := \bigl(\|v\|_{L^2(\dom)}^2 + \|v'\|_{L^2(\dom)}^2\bigr)^{1/2}.
\end{equation}
For homogeneous Dirichlet conditions, we set
\begin{equation}
V := H_0^1(\dom) := \bigl\{v \in H^1(\dom) : v(0) = v(\ell) = 0\bigr\}.
\end{equation}
By the Poincar\'e inequality, for $v \in H_0^1(\dom)$,
\begin{equation}\label{eq:poincare}
\|v\|_{L^2(\dom)} \leq C_P\,\|v'\|_{L^2(\dom)}, \qquad C_P = \frac{\ell}{\pi}.
\end{equation}

For non-homogeneous data $g_0, g_\ell$, we choose any lift $w_g \in H^1(\dom)$ with $w_g(0)=g_0$, $w_g(\ell)=g_\ell$ and set $u = w_g + w$ where $w \in H_0^1(\dom)$. In the following we assume homogeneous conditions for clarity; the non-homogeneous case follows by a standard lifting argument.

\begin{definition}[Bilinear form and linear functional]
\label{def:bilinear}
Define $a : H^1(\dom) \times H^1(\dom) \to \R$ and $L : H^1(\dom) \to \R$ by
\begin{align}
a(w,v) &:= \int_0^\ell \kappa\,w'v'\dd x + \int_0^\ell \beta\,w'v\dd x + \int_0^\ell \sigma\,wv\dd x,\label{eq:bilinear}\\
L(v) &:= \int_0^\ell f\,v\dd x.\label{eq:linear}
\end{align}
\end{definition}

The \textbf{weak formulation} of \eqref{eq:PDE} is: find $u \in V$ such that
\begin{equation}\label{eq:weak}
a(u,v) = L(v) \qquad \forall v \in V.
\end{equation}

\begin{proposition}[Continuity and coercivity]
\label{prop:continuity_coercivity}
The bilinear form $a(\cdot,\cdot)$ is continuous on $V$ with constant
\begin{equation}\label{eq:M}
M := \kappa + C_P\,\|\beta\|_{L^\infty(\dom)} + C_P^2\,\|\sigma\|_{L^\infty(\dom)},
\end{equation}
that is, $|a(w,v)| \leq M\,\|w\|_V\,\|v\|_V$ for all $w,v \in V$, where $\|v\|_V := \|v'\|_{L^2(\dom)}$.

If $\sigma \geq 0$ a.e.\ and $\beta$ is constant, then $a$ is coercive on $V$ with constant
\begin{equation}\label{eq:alpha}
\alpha = \kappa > 0,
\end{equation}
that is, $a(v,v) \geq \alpha\,\|v\|_V^2$ for all $v \in V$.
\end{proposition}

\begin{proof}
\textbf{Continuity.} By the Cauchy--Schwarz inequality and \eqref{eq:poincare}:
\begin{align}
|a(w,v)| &\leq \kappa\|w'\|_{L^2}\|v'\|_{L^2} + \|\beta\|_{L^\infty}\|w'\|_{L^2}\|v\|_{L^2} + \|\sigma\|_{L^\infty}\|w\|_{L^2}\|v\|_{L^2}\\
&\leq \bigl(\kappa + C_P\|\beta\|_{L^\infty} + C_P^2\|\sigma\|_{L^\infty}\bigr)\|w'\|_{L^2}\|v'\|_{L^2}.
\end{align}

\textbf{Coercivity.} For constant $\beta$ and $\sigma \geq 0$:
\begin{equation}
a(v,v) = \kappa\|v'\|_{L^2}^2 + \int_0^\ell \beta\,v'v\dd x + \sigma\|v\|_{L^2}^2.
\end{equation}
The advection term integrates to zero by integration by parts:
\begin{equation}
\int_0^\ell \beta\,v'v\dd x = \frac{\beta}{2}\bigl[v^2\bigr]_0^\ell = \frac{\beta}{2}\bigl(v(\ell)^2 - v(0)^2\bigr) = 0,
\end{equation}
since $v(0) = v(\ell) = 0$. Therefore $a(v,v) = \kappa\|v'\|_{L^2}^2 + \sigma\|v\|_{L^2}^2 \geq \kappa\,\|v\|_V^2$.
\end{proof}

By the Lax--Milgram theorem \cite{brezis2010functional}, the weak problem \eqref{eq:weak} has a unique solution $u \in V$.

\subsection{The classical Galerkin approximation}

Let $\Th = \{K_j = (x_j, x_{j+1}) : j = 0,\ldots,N-1\}$ be a partition of $\dom$ into intervals with element sizes $h_j = x_{j+1} - x_j$ and global mesh size $h = \max_j h_j$. Define the piecewise linear finite element space:
\begin{equation}\label{eq:VhP}
\Vp := \bigl\{v_h \in V : v_h\big|_K \in \mathbb{P}_1(K) \text{ for all } K \in \Th\bigr\},
\end{equation}
where $\mathbb{P}_1(K)$ denotes the space of polynomials of degree $\leq 1$ on $K$. The dimension of $\Vp$ is $N-1$ (the interior nodes).

The \textbf{classical Galerkin method} seeks $u_h \in \Vp$ satisfying
\begin{equation}\label{eq:Galerkin}
a(u_h, v_h) = L(v_h) \qquad \forall v_h \in \Vp.
\end{equation}

The error of the Galerkin approximation is governed by \textbf{C\'ea's lemma} \cite{brenner2008mathematical}:

\begin{theorem}[C\'ea's lemma]
\label{thm:cea}
Let $u \in V$ solve \eqref{eq:weak} and $u_h \in \Vp$ solve \eqref{eq:Galerkin}. Then
\begin{equation}\label{eq:cea}
\|u - u_h\|_V \leq \frac{M}{\alpha}\,\inf_{v_h \in \Vp}\|u - v_h\|_V.
\end{equation}
For the problem \eqref{eq:PDE} with constant $\beta$ and $\sigma \geq 0$, this becomes
\begin{equation}\label{eq:cea_explicit}
\|u - u_h\|_V \leq \left(1 + C_P\,\frac{\|\beta\|_{L^\infty}}{\kappa} + C_P^2\,\frac{\|\sigma\|_{L^\infty}}{\kappa}\right)\inf_{v_h \in \Vp}\|u - v_h\|_V.
\end{equation}
\end{theorem}

\subsection{Failure in the singularly perturbed regime}

The ratio $M/\alpha$ in \eqref{eq:cea_explicit} is precisely the \emph{global condition number} of the bilinear form:
\begin{equation}\label{eq:condition}
\frac{M}{\alpha} = 1 + C_P\,\frac{\|\beta\|_{L^\infty}}{\kappa} + C_P^2\,\frac{\|\sigma\|_{L^\infty}}{\kappa}.
\end{equation}

For a standard second-order approximation, the best approximation error satisfies $\inf_{v_h}\|u - v_h\|_V \leq C\,h\,|u|_{H^2(\dom)}$. Combining with \eqref{eq:cea} gives:
\begin{equation}\label{eq:error_Galerkin}
\|u - u_h\|_V \leq \frac{M}{\alpha}\,C\,h\,|u|_{H^2(\dom)}.
\end{equation}

The critical difficulty is twofold:
\begin{enumerate}[label=(\roman*)]
    \item The \textbf{condition number} $M/\alpha$ grows linearly with $\|\beta\|/\kappa$ and quadratically with $\|\sigma\|/\kappa$, so the constant in \eqref{eq:error_Galerkin} deteriorates as the problem becomes singularly perturbed.
    \item The \textbf{regularity} of $u$ degrades simultaneously: when $\Pe \gg 1$, the solution develops an exponential boundary layer of width $\delta \sim \kappa/|\beta|$ at the outflow boundary, and $|u|_{H^2(\dom)}$ grows like $O((|\beta|/\kappa)^2)$. The product $M/\alpha \cdot |u|_{H^2}$ therefore grows as $O(\Pe^3)$.
\end{enumerate}

In practice, the classical Galerkin approximation produces \emph{spurious oscillations} near layers when $\Pe_K > 1$, which pollute the global solution. This motivates enrichment of the finite element space with local subscale functions.

\subsection{Outline}

The remainder of this article is organized as follows. \Cref{sec:RFB} presents the complete theory of residual-free bubbles: the enriched space, local bubble problems, dimension reduction via residual-mode bases, static condensation, and the error theorem for approximate bubbles. \Cref{sec:KAN} describes the KAN architecture, its mathematical foundation in the Kolmogorov--Arnold representation theorem, and its B-spline parametrization. \Cref{sec:learning} explains how the KAN is trained as a surrogate for the exact bubbles. \Cref{sec:theory} provides rigorous error estimates with full proofs. \Cref{sec:numerics} outlines the numerical experiments, and \Cref{sec:conclusion} concludes.

% ════════════════════════════════════════════════════════════════════════
\section{Residual-Free Bubble Functions}
\label{sec:RFB}

\subsection{Enriched space and the local bubble problem}

We enrich the coarse space $\Vp$ by adding, on each element $K \in \Th$, a finite-dimensional subspace $B^K \subset H_0^1(K)$ whose functions are extended by zero outside $K$. Define the bubble space and the enriched space:
\begin{equation}\label{eq:BK_def}
B := \bigoplus_{K \in \Th} B^K, \qquad \Vh := \Vp \oplus B.
\end{equation}
Any function $v_h \in \Vh$ splits uniquely as
\begin{equation}
v_h = v_P + \sum_{K \in \Th} v_{B,K}, \qquad v_P \in \Vp,\quad v_{B,K} \in B^K.
\end{equation}

The Galerkin approximation on the enriched space reads: find $u_h = u_P + u_B \in \Vh$ such that
\begin{equation}\label{eq:enriched_split}
a(u_h, v_P) = L(v_P) \quad \forall v_P \in \Vp,
\qquad
a(u_h, v_{B,K})_K = L_K(v_{B,K}) \quad \forall v_{B,K} \in B^K,\;\forall K \in \Th,
\end{equation}
where the subscript $K$ restricts integrals to element $K$.

Since $u_h\big|_K = u_P\big|_K + u_{B,K}$, the second equation becomes, for each $K$:
\begin{equation}\label{eq:RFB_residual}
a_K(u_{B,K}, v_{B,K}) = -\bigl(a_K(u_P, \cdot) - L_K(\cdot)\bigr)(v_{B,K}) \qquad \forall v_{B,K} \in B^K.
\end{equation}
The right-hand side is the \emph{element residual} of the coarse function.

\subsection{The residual-free condition}
\label{subsec:RFB_condition}

The residual-free bubble approach \cite{babuska1994special,brezzi1994further} requires that \eqref{eq:RFB_residual} hold for \emph{every} test function in $H_0^1(K)$, not just in $B^K$. That is, $u_{B,K}$ satisfies:
\begin{equation}\label{eq:RFB_exact}
a_K(u_{B,K}, v) = -\bigl(a_K(u_P, \cdot) - L_K(\cdot)\bigr)(v) \qquad \forall v \in H_0^1(K).
\end{equation}

By the Lax--Milgram theorem, this problem has a unique solution in $H_0^1(K)$ for each fixed $u_P\big|_K$. This defines an affine operator:
\begin{equation}
\calT_K : \mathbb{P}_1(K) \to H_0^1(K), \qquad u_P\big|_K \mapsto u_{B,K}.
\end{equation}
The residual-free bubble space is then
\begin{equation}
B^K := \calT_K\bigl(\mathbb{P}_1(K)\bigr) \subset H_0^1(K).
\end{equation}

The key property of the RFB enrichment is that the coarse function $u_P$ and the bubble correction $u_{B,K}$ together satisfy the original PDE exactly in the interior of each element:
\begin{equation}
\calL(u_P + u_{B,K}) = f \quad \text{in } K, \qquad u_{B,K} = 0 \quad \text{on } \partial K,
\end{equation}
where $\calL$ denotes the differential operator in \eqref{eq:PDE}. The residual of the coarse function is exactly cancelled by the bubble correction.

\subsection{Residual-mode basis for $B^K$}
\label{subsec:residual_modes}

In one dimension with $\mathbb{P}_1$ coarse space, each element $K$ contributes two nodal basis functions $\phi_{1,K}$, $\phi_{2,K}$. Writing $u_P\big|_K = U_1^K \phi_{1,K} + U_2^K \phi_{2,K}$, the affine nature of \eqref{eq:RFB_exact} gives
\begin{equation}
u_{B,K} = U_1^K\,\calT_K(\phi_{1,K}) + U_2^K\,\calT_K(\phi_{2,K}) + \calT_K(0; f),
\end{equation}
where $\calT_K(0;f)$ is the bubble generated by the source alone. Thus $\dim B^K \leq 3$ in general.

On the reference element $\widehat{K} = (0,1)$ with coordinate $\xi = (x - x_j)/h_K$, the differential operator becomes
\begin{equation}\label{eq:L_ref}
\calL_\xi\, w := -\frac{\kappa}{h_K^2}\,w'' + \frac{\beta}{h_K}\,w' + \sigma\,w,
\end{equation}
and the strong form of the bubble problem is
\begin{equation}\label{eq:local_strong}
-\frac{\kappa}{h_K^2}\,b'' + \frac{\beta}{h_K}\,b' + \sigma\,b = r(\xi), \qquad b(0) = b(1) = 0,
\end{equation}
where $r(\xi)$ is the element residual restricted to $K$.

The explicit expressions for the three residual contributions are:
\begin{align}
-\calL_\xi\,\phi_{1,K} &= \frac{\beta}{h_K} - \sigma(1-\xi) = \Bigl(\frac{\beta}{h_K} - \sigma\Bigr) + \sigma\,\xi,\\
-\calL_\xi\,\phi_{2,K} &= -\frac{\beta}{h_K} - \sigma\,\xi,\\
f(\xi) &\equiv 1 \quad \text{(for a constant source)}.
\end{align}
Every right-hand side is an affine function of $\xi$. Therefore, it suffices to learn only two \emph{residual-mode bubbles}:
\begin{equation}\label{eq:residual_modes}
\widehat{b}(\xi) := \calL_\xi^{-1}(1), \qquad \widetilde{b}(\xi) := \calL_\xi^{-1}(\xi),
\end{equation}
with homogeneous Dirichlet conditions on $(0,1)$. Any bubble arising from an affine right-hand side is a linear combination of $\widehat{b}$ and $\widetilde{b}$. In particular:
\begin{equation}\label{eq:bk_components}
b_{1,K} = \Bigl(\frac{\beta}{h_K} - \sigma\Bigr)\widehat{b} + \sigma\,\widetilde{b}, \qquad
b_{2,K} = -\frac{\beta}{h_K}\,\widehat{b} - \sigma\,\widetilde{b}, \qquad
b_K^f = \widehat{b}.
\end{equation}
Thus the bubble space has dimension at most 2:
\begin{equation}
B^K = \operatorname{span}\{\widehat{b}(\xi),\;\widetilde{b}(\xi)\}.
\end{equation}

\begin{remark}[Normalization]
\label{rem:normalize}
Each residual-mode bubble is normalized by its value at the midpoint:
\begin{equation}
\widehat{b}_{\mathrm{norm}}(\xi) := \frac{\widehat{b}(\xi)}{\widehat{b}(1/2)}, \qquad
\widetilde{b}_{\mathrm{norm}}(\xi) := \frac{\widetilde{b}(\xi)}{\widetilde{b}(1/2)},
\end{equation}
yielding $\widehat{b}_{\mathrm{norm}}(0) = \widehat{b}_{\mathrm{norm}}(1) = 0$ and $\widehat{b}_{\mathrm{norm}}(1/2) = 1$. This fixes the arbitrary multiplicative scale.
\end{remark}

\begin{remark}[Maximum principle]
\label{rem:max_principle}
The operator $\calL_\xi$ satisfies the classical one-dimensional maximum principle when $\kappa > 0$ and $\sigma \geq 0$. Consequently, if $r(\xi) \geq 0$ on $K$, then $b = \calL_\xi^{-1}(r) \geq 0$ everywhere on $[0,1]$. For $\beta > 0$ and $\sigma = 0$, both $\widehat{b} = \calL_\xi^{-1}(1) \geq 0$ and $\widetilde{b} = \calL_\xi^{-1}(\xi) \geq 0$.
\end{remark}

\subsection{Static condensation}
\label{sec:condensation}

Let the two residual-mode bubbles on element $K$ be collected as $\bm{b}_K = (\widehat{b}_K, \widetilde{b}_K)^T$. The enriched approximation on $K$ is
\begin{equation}
u_h\big|_K = U_1^K \phi_{1,K} + U_2^K \phi_{2,K} + U_{\widehat{b}}^K \widehat{b}_K + U_{\widetilde{b}}^K \widetilde{b}_K.
\end{equation}
Collecting coarse and bubble unknowns into vectors
\begin{equation}
U_L^K = \begin{bmatrix} U_1^K \\ U_2^K \end{bmatrix} \in \R^2, \qquad
U_b^K = \begin{bmatrix} U_{\widehat{b}}^K \\ U_{\widetilde{b}}^K \end{bmatrix} \in \R^2,
\end{equation}
the element system becomes
\begin{equation}\label{eq:element_system}
\begin{bmatrix}
A_{LL}^K & A_{Lb}^K \\
A_{bL}^K & A_{bb}^K
\end{bmatrix}
\begin{bmatrix}
U_L^K \\ U_b^K
\end{bmatrix}
=
\begin{bmatrix}
F_L^K \\ F_b^K
\end{bmatrix},
\end{equation}
where
\begin{align}
(A_{LL}^K)_{ij} &= a_K(\phi_j^K, \phi_i^K), & (A_{Lb}^K)_{i\mu} &= a_K(b_{K,\mu}, \phi_i^K), \\
(A_{bL}^K)_{\mu j} &= a_K(\phi_j^K, b_{K,\mu}), & (A_{bb}^K)_{\mu\nu} &= a_K(b_{K,\nu}, b_{K,\mu}), \\
(F_L^K)_i &= L_K(\phi_i^K), & (F_b^K)_\mu &= L_K(b_{K,\mu}).
\end{align}
for $i,j \in \{1,2\}$ and $\mu,\nu \in \{\widehat{b}, \widetilde{b}\}$.

Solving the second block row for the bubble coefficients:
\begin{equation}\label{eq:bubble_recovery}
U_b^K = (A_{bb}^K)^{-1}\bigl(F_b^K - A_{bL}^K\,U_L^K\bigr).
\end{equation}
Substituting into the first block row yields the \textbf{condensed system}:
\begin{equation}\label{eq:condensed}
A_{\mathrm{cond}}^K\,U_L^K = F_{\mathrm{cond}}^K,
\end{equation}
with the Schur complement
\begin{equation}\label{eq:schur}
A_{\mathrm{cond}}^K = A_{LL}^K - A_{Lb}^K\,(A_{bb}^K)^{-1}\,A_{bL}^K, \qquad
F_{\mathrm{cond}}^K = F_L^K - A_{Lb}^K\,(A_{bb}^K)^{-1}\,F_b^K.
\end{equation}

Only the condensed matrices are assembled into the global system, which therefore has exactly $\dim \Vp$ unknowns---the same as the classical P1 method. After the global solve, bubble coefficients are recovered elementwise via \eqref{eq:bubble_recovery}. The Schur complement $-A_{Lb}^K(A_{bb}^K)^{-1}A_{bL}^K$ adds a dense $2 \times 2$ correction per element without introducing new global couplings; the sparsity pattern remains tridiagonal.

% ════════════════════════════════════════════════════════════════════════
\bibliographystyle{plain}
\bibliography{references}

\end{document}
